% Figure from tools/make_paper_figures.py.
\section*{Discussion}
\label{discussion}

The clearest result is a disagreement between the two kinds of measurement.
Qwen3 is closest to the patient by speaker embedding, so close that ECAPA-TDNN
cannot separate it from a second genuine recording, and it is the system raters
find least natural and detect most reliably. Voxtral is the mirror image
(Figure~\ref{fig:sim_vs_nmos}). Both perceptual measures agree with each other
and contradict the objective one; a selection by speaker similarity alone would
have picked the worst system for this purpose.

\begin{figure}[H]
    \centering
    \includegraphics[width=0.45\linewidth]{fig_sim_vs_nmos.pdf}
    \caption{Speaker similarity against naturalness. The system closest to the
    patient by embedding is the one raters find least natural.}
    \label{fig:sim_vs_nmos}
\end{figure}

The disagreement concerns similarity, not objective measurement as such.
Response time cannot be obtained from listeners, and it is the practical
constraint: every service renders faster than real time, yet median end-to-end
latency ran from $5.7$ to $13.8$\,s per utterance because the texts are long. A
conversational application would have to synthesise in shorter units.

The clones do not pass as the patient. S2 Pro was identified in three of four
pairs, the patient herself identified her voice in six of six. Legacy material
yields a voice listeners accept as natural, not one that survives a direct
comparison with her own; for an assistive device, which speaks for the patient
rather than deceiving anyone, that is the relevant standard.

The reversed familiarity effect cannot be interpreted. No unfamiliar rater was
older than $42$, the familiar group has a median age of $51$, and an
age-related loss of high-frequency hearing would blunt exactly the brightness
cue that separates genuine from cloned excerpts; hearing was not screened.
Beyond this, the study covers one speaker, the fooling rates are a lower bound
because pause structure, brightness and duration remain usable cues, and the
surviving texts are those Voxtral rendered, which flatters it.

The same property that makes this result useful is also a risk. The material
was not recorded for cloning; it is the kind of audio most people have sent to
others. That it suffices for a voice which unfamiliar listeners mistake for the
original in up to three of ten side-by-side comparisons means the barrier to
voice fraud is no longer expertise or studio recordings but access to someone's
messages, and a fraudulent call never has to survive a side-by-side. This is no
reason to withhold the method from patients, but it is a reason to treat legacy
audio as sensitive data.

For the patient, the primary result stands: two of the three systems produced a
voice that $30$ listeners rated as natural as her own recordings, from material
never intended for this purpose. That the two can still be told apart in a
direct comparison limits the clone as a replica, not as a voice.
